\section{Weighted annular averages and the peripheral spectral measure}
\label{sec:weighted-spectral-averages}
The direct orbit estimate has two consequences which require no
variation bound for the final test vector: cancellation under an
unchanged event, and control of the spectral measure of the constant
vector on every peripheral arc. We also identify the small-frequency
limit of that spectral measure. These are statements about the exact
unitary return dynamics, at the level of averages or of its specified
cyclic vector. They are not operator-norm power bounds.

\subsection{Weights and unchanged events}
For a complex $a\in\mathcal H_R$ define
\begin{equation}\label{eq:averaged-weighted-characteristic}
 C^a_{j,R}(z)=\int a(x)e^{-iz\cdot U_{j,R}(x)}\dd\nu_R^*(x),
       \qquad j\ge0.
\end{equation}
The weight is not normalized. Let $\mathcal A_n$ be the physical
annulus in \eqref{eq:direct-annulus}.

\begin{theorem}[Mean-square cancellation with arbitrary square-integrable weights]
\label{thm:weighted-annular-averages}
For $n$ sufficiently large, $T\ge n$, $N\ge0$, and every
$z\in\mathcal A_n$,
\begin{equation}\label{eq:weighted-annular-averages}
 \frac1T\sum_{j=0}^{T-1}|C^a_{N+j,R}(z)|^2
                    \le Cn^{-4/495}\|a\|_2^2.
\end{equation}
If $A$ is any measurable event with $p=\nu_R^*(A)>0$, then
\begin{equation}\label{eq:unchanged-event-average}
 \frac1T\sum_{j=0}^{T-1}
 \left|\mathbb E_{\nu_R^*}
     [e^{-iz\cdot U_{N+j,R}}\mid A]\right|^2
                    \le \frac C p n^{-4/495}.
\end{equation}
The same event $A$ occurs at every averaging index. The constants are
uniform in $R,N,T,z,a,A$ subject to these conditions.
\end{theorem}
\begin{proof}
Take $b=\overline a$ in
\eqref{eq:direct-moving-analysis} and use
Corollary~\ref{cor:direct-annular-averages}. For the conditional
statement take $a=\one_A/p$, whose squared $L^2$ norm is exactly
$1/p$. This proves both estimates with their actual normalizations.
No smoothing, count restriction, subanalytic assumption or factorization
of the event probability is used.
\end{proof}

For example, if $p\ge p_0 n^{-\beta}$ with $\beta<4/495$, the
average in \eqref{eq:unchanged-event-average} tends to zero uniformly.
More quantitatively, for every $\eta>0$, the fraction of indices
$j<T$ at which the displayed conditional transform exceeds $\eta$
in modulus is at most
\begin{equation}\label{eq:exceptional-return-fraction}
 C\eta^{-2}p^{-1}n^{-4/495}.
\end{equation}
This follows by applying the elementary counting inequality to the
sum of squares. The exceptional set may depend on the frequency and
the event. No claim is made that one specified return count avoids it.

\begin{corollary}[A fixed actual mark and a moving terminal mark]
\label{cor:averaged-actual-marks}
The bound \eqref{eq:weighted-annular-averages} also holds for
\[
 \int a((F_R^*)^k x)e^{-iz\cdot U_{N+j,R}(x)}\dd\nu_R^*(x)
\]
when $0\le k\le N$ is fixed throughout the averaging window. It holds
as well when the weight is $a((F_R^*)^{N+j}x)$, evaluated at the actual
terminal return at each index.
\end{corollary}
\begin{proof}
Write $\psi_l=\mathcal U_z^l e$, including negative integers.
The weighted-composition identity for $\mathcal U_z^{-k}$ is
\[
 \psi_{l-k}(y)=\psi_{-k}(y)\psi_l((F_R^*)^{-k}y).
\]
After changing variables $y=(F_R^*)^kx$, the fixed-mark integral is
$\langle\overline a\,\psi_{-k},\psi_{N+j-k}\rangle_R$.
Its test vector has norm $\|a\|_2$ and is independent of $j$.
Apply \eqref{eq:direct-moving-analysis}.

For the terminal mark, change variables $y=(F_R^*)^{N+j}x$.
The phase becomes $\overline{\psi_{-(N+j)}(y)}$. The conjugate of
the resulting integral is $\langle a,\psi_{-(N+j)}\rangle_R$.
Apply the same theorem to the inverse unitary, whose absolute lag
correlations are unchanged. These changes of variables use the actual
return map and do not replace a random observation by a deterministic
collision-clock observation.
\end{proof}

\begin{corollary}[A normalized annular integral]
\label{cor:normalized-annular-integral}
For the original law, and uniformly in $R,N\ge0,T\ge n$,
\begin{equation}\label{eq:normalized-annular-integral}
 \frac1{|\mathcal A_n|}\int_{\mathcal A_n}
     \frac1T\sum_{j=0}^{T-1}|c_{N+j,R}(z)|^2\dd z
                                      \le Cn^{-4/495}.
\end{equation}
The same conclusion holds for an event as in
\eqref{eq:unchanged-event-average}, with the factor $p^{-1}$.
\end{corollary}
\begin{proof}
Integrate the uniform bounds in
Theorem~\ref{thm:weighted-annular-averages}, and divide by the volume
of this nonempty annulus. Finite sums and bounded integrands justify
interchanging the integrals.
\end{proof}

The normalization in \eqref{eq:normalized-annular-integral} is stated
explicitly. The raw local theorem instead requires an \emph{unnormalized}
complementary integral at a specified return count, with the four-dimensional
factor $n^2$. These requirements are not interchangeable. Indeed
$|\mathcal A_n|\asymp n^{-8/5}$; Cauchy--Schwarz applied to the new
mean-square bound would give at best
$n^{2/5-2/495}$ after that raw normalization. It therefore supplies
no vanishing raw error by itself. Likewise a fixed event across an
averaging window is different from a separately chosen exact physical
observation event for every member of that window.

\subsection{Every peripheral arc for the constant vector}
Let $\varsigma_{R,z}$ be the probability on $[-\pi,\pi)$ associated
by the unitary spectral theorem with $e=\one$ and $\mathcal U_z$.
Thus
\begin{equation}\label{eq:cyclic-spectral-measure}
 c_{j,R}(z)=\int e^{ij\theta}\dd\varsigma_{R,z}(\theta),
                         \qquad j\in\Z.
\end{equation}
This is a spectral measure of the return operator, not the time density
of the physical record. The distinction is relevant to both of the
following statements.

\begin{theorem}[Small peripheral arcs and exterior cyclic resolvents]
\label{thm:cyclic-peripheral-bounds}
For $0<|z|=r\le r_0$, $m=\lfloor r^{-200/99}\rfloor$, and all
$\xi\in\R$,
\begin{align}
 \varsigma_{R,z}\{\theta:
   \operatorname{dist}(\theta-\xi,2\pi\Z)\le m^{-1}\}
                                      &\le Cr^{2/99},
                              \label{eq:cyclic-arc-bound}\\
 \left\|(1-s)(I-se^{-i\xi}\mathcal U_z)^{-1}e\right\|_2^2
                                      &\le Cr^{2/99}
                              \label{eq:cyclic-Abel-bound}
\end{align}
whenever $0<s<1$ and $(1-s)m\le1$. All constants are uniform in $R$.
The inverse in \eqref{eq:cyclic-Abel-bound} is an ordinary exterior
unitary resolvent applied to one specified vector. No inverse of
$e^{i\xi}-\mathcal U_z$ is asserted on the unit circle.
\end{theorem}
\begin{proof}
The squared norm of
$m^{-1}\sum_{j<m}e^{-ij\xi}\mathcal U_z^je$ is the integral against
$\varsigma_{R,z}$ of
\[
 \left|\frac1m\sum_{j=0}^{m-1}e^{ij(\theta-\xi)}\right|^2.
\]
On the arc of radius $1/m$ this kernel is bounded below by a positive
absolute constant. For example, use
$|\sin(m t/2)|/(m|\sin(t/2)|)\ge2/\pi$ for $|t|\le1/m$,
with its continuous value at zero. Equation~\eqref{eq:orbit-synthesis}
and Proposition~\ref{prop:short-lag-row-budget} bound the total integral
by $Cr^{2/99}$. This proves \eqref{eq:cyclic-arc-bound} simultaneously
for all centers, with no choice of peripheral representatives.

For the second assertion expand the resolvent in its norm-convergent
Neumann series. In its block of indices $km,\ldots,km+m-1$,
\eqref{eq:orbit-synthesis} gives the norm bound
\[
 s^{km}\sqrt{\mathcal B_m\sum_{j=0}^{m-1}s^{2j}}
                         \le s^{km}\sqrt{m\mathcal B_m}.
\]
Sum the blocks and multiply by $1-s$. Since
$s^m\le e^{-m(1-s)}$ and
$1-e^{-x}\ge(1-e^{-1})x$ for $0\le x\le1$, the result is at most
$C\sqrt{\mathcal B_m/m}$. Squaring and applying
\eqref{eq:short-lag-row-budget} proves \eqref{eq:cyclic-Abel-bound}.
The proof controls no operator norm on arbitrary input vectors.
\end{proof}

\subsection{A Cauchy scaling law for the exact spectral measure}
For a unit direction $u\in\R^4$ put
\[
 d_{R,u}=u^{\mathsf T}D_Ru,\qquad a_{R,u}=d_{R,u}/2.
\]
Uniform ellipticity bounds $a_{R,u}$ above and below by positive
constants. Let $\mathsf C_a$ be the Cauchy probability on $\R$ with
density $a/[\pi(a^2+x^2)]$ and characteristic function $e^{-a|t|}$.
The principal angle in \eqref{eq:cyclic-spectral-measure} is used
in the rescaling below.

\begin{theorem}[Spectral Cauchy limit at physical frequency zero]
\label{thm:spectral-Cauchy-limit}
Let $\widetilde\varsigma_{R,r,u}$ be the pushforward of
$\varsigma_{R,ru}$ under $\theta\mapsto\theta/r^2$.
Then, as $r\downarrow0$,
\begin{equation}\label{eq:spectral-Cauchy-limit}
 \sup_{R\in I,\ |u|=1}
 d_{\rm BL}\bigl(\widetilde\varsigma_{R,r,u},
                         \mathsf C_{a_{R,u}}\bigr)\longrightarrow0.
\end{equation}
This describes a spectral probability. It does not claim that the
physical return record has a Cauchy limit; its Gaussian limit is
unchanged.
\end{theorem}
\begin{proof}
The one-return second moment is uniformly bounded by
Theorem~\ref{thm:exponential-return}. Using $1-\cos x\le x^2/2$,
\[
 1-\operatorname{Re}c_{1,R}(ru)
      \le Cr^2.
\]
For $|\theta|\le\pi$ one has
$\theta^2\le(\pi^2/2)(1-\cos\theta)$. Therefore
\begin{equation}\label{eq:principal-angle-rounding}
 \int|\theta|\dd\varsigma_{R,ru}(\theta)\le Cr.
\end{equation}
Fix $t>0$ and put $j=\lfloor t/r^2\rfloor$. The difference between
the characteristic function of $\widetilde\varsigma_{R,r,u}$ at $t$
and $c_{j,R}(ru)$ is at most $Cr$, by
\eqref{eq:principal-angle-rounding} and $|t/r^2-j|<1$.
For sufficiently small $r$, $j\ge2$, and Theorem A applied at
rescaled frequency $-r\sqrt j\,u$ gives
\[
 c_{j,R}(ru)=\exp(-jr^2d_{R,u}/2)+
               O_t\bigl(j^{-1/26}\sqrt{\log(2+j)}\bigr).
\]
The rescaled frequencies remain in a fixed ball for fixed $t$.
The main term tends uniformly to $e^{-td_{R,u}/2}$. Negative $t$
uses $c_{-j}=\overline{c_j}$, and $t=0$ is exact. This proves
pointwise characteristic convergence, uniformly in the parameters,
to the asserted Cauchy family.

To justify the uniform bounded-Lipschitz conclusion, take any sequence
$r_l\downarrow0$, $R_l\in I$, $|u_l|=1$. Compactness gives a
subsequence $R_l\to R_*$, $u_l\to u_*$. Continuity of $D_R$ and the
preceding characteristic estimates identify its limit as
$\mathsf C_{a_{R_*,u_*}}$ by the characteristic-function continuity
theorem. The target Cauchy laws converge as well: multiply one standard
Cauchy random variable by $a_{R_l,u_l}$. Thus a sequence violating
\eqref{eq:spectral-Cauchy-limit} cannot exist. No absolute continuity
of $\varsigma_{R,z}$ for a fixed nonzero $z$ is presumed or deduced.
\end{proof}

\paragraph{Place in the raw inversion problem.}
Theorems~\ref{thm:direct-moving-averages},
\ref{thm:weighted-annular-averages} and
\ref{thm:cyclic-peripheral-bounds} bypass the incompatible grid scales
for averaged physical coefficients and for a specified cyclic
resolvent vector. They cover every peripheral angle and allow
square-integrable weights without a BV reconstruction hypothesis.
Theorem~\ref{thm:spectral-Cauchy-limit} identifies the concentration
scale of the same spectral probability. None of these statements
replaces the fixed-return complementary integral in
Theorem~\ref{thm:finite-extraction-central-inversion}, the long-time growth of
$A_2(n,L_n,R,w)$, or the local extracted-edge correction. The exact
physical-event replacement remains a different assertion from the
unchanged-event average proved above.
